% !TEX program = xelatex

\documentclass{resume}
\usepackage{setspace}

\begin{document}
\pagenumbering{gobble} % suppress displaying page number

\name{Wenxuan Wu}

\basicInfo{
  \email{wwx0306@foxmail.com} \textperiodcentered\
  \phone{(+86) 191-159-46187} \textperiodcentered\
  \github[Wenxuan Wu]{https://github.com/wuwenxux}}

\section{\faRocket\ Strengths}
\begin{itemize}[parsep=0.5ex]
	\item Extensive server/embedded Linux development experience with deep expertise in kernel network stack
	\item RISC-V ISA: DPDK/SPDK driver porting, QEMU device emulation, Ceph distributed storage deployment and optimization
	\item DPDK PMD driver development (Intel ixgbe/i40e/ice); custom transport-layer protocol design; open-source contributions to Kernel and DPDK
\end{itemize}

\section{\faBriefcase\ Work Experience}
\datedsubsection{\textbf{Rivosemi}, Shenzhen}{Jul. 2025 -- Present}
\role{Low-Level Software Engineer}{RISC-V C Ceph DPDK SPDK QEMU Shell}
\begin{onehalfspacing}
Ceph Distributed Storage Deployment and DPDK/SPDK Porting on RISC-V
\begin{itemize}
	\item Compiled Ceph 19.2.3 from source on RISC-V with DPDK/SPDK, fixed 5 RISC-V specific compilation errors, verified ceph-mon/osd linking with DPDK
	\item Built ~8,500-line Shell deployment toolkit (7 scripts) for one-click 3-node Ceph cluster with auto-recovery from 14 failure types
	\item Validated DPDK/SPDK on RISC-V; delivered DPDK implementation roadmap and RVA23 ISA performance analysis for SPDK/RDMA critical paths
\end{itemize}
\end{onehalfspacing}

\begin{onehalfspacing}
QEMU Emulation Platform Based on Proprietary RISC-V SoC
\begin{itemize}
	\item Built full-system QEMU emulation for proprietary RISC-V SoC enabling pre-silicon OS boot, driver, storage/network I/O validation; JSON-driven device framework covering 6 device types
	\item Designed fault injection framework (28 types) and phased performance injection engine; refactored SoC codebase from single file to 8 independent modules
\end{itemize}
\end{onehalfspacing}

\datedsubsection{\textbf{Autel Robotics Co., Ltd.}, Shenzhen}{Feb. 2025 -- Apr. 2025}
\role{Senior Network Development Engineer}{ARM64 C OpenSSL GmSSL Network Security}
\begin{onehalfspacing}
GmSSL Encryption Implementation for UAV Communication Networks
\begin{itemize}
	\item Implemented GmSSL encryption for power-industry requirements, establishing full-link encrypted channel across UAV, remote controller, and base station
	\item Analyzed ChannelManager networking protocol interactions, authored end-to-end encryption testcases
\end{itemize}
\end{onehalfspacing}


\datedsubsection{\textbf{TD Tech Ltd.}, Chengdu}{Jul. 2023 -- Feb. 2025}
\role{Senior Network Development Engineer}{C gdb Kernel Network Stack HW/SW Co-Design}
\begin{onehalfspacing}
Kernel Network Stack Maintenance and Performance Optimization for Module Terminals
\begin{itemize}
	\item Led network protocol stack performance analysis and bottleneck identification across mips, arm, and x86 CPU architectures
	\item Leveraged hardware offloading capabilities (checksum, segmentation offload) to reduce kernel stack processing overhead and improve data plane throughput
\end{itemize}
\end{onehalfspacing}

\datedsubsection{\textbf{Seawave Technology}, Chengdu}{Jul. 2022 -- Aug. 2023}
\role{DPDK Network Protocol Stack Engineer}{C DPDK Custom Transport Layer Protocol Design}
\begin{onehalfspacing}
Custom Transport-Layer Protocol Design and Implementation
\begin{itemize}
	\item Independently designed and implemented transport-layer protocol mechanisms: out-of-order reordering and NACK-based retransmission for reliable delivery without TCP
	\item Independently developed network traffic generation tool supporting overlay, multipath, and covert communication scenarios
	\item Served as cache module Maintainer; completed asset documentation (design docs, usage guides, bug fixes) and code review
\end{itemize}
\end{onehalfspacing}

\datedsubsection{\textbf{Thundersoft Co., Ltd.}, Chengdu}{Nov. 2021 -- Jul. 2022}
\role{DPDK Driver Development Engineer}{C DPDK Intel NIC Performance Optimization}
\begin{onehalfspacing}
Intel Ethernet PMD Driver Development and Maintenance
\begin{itemize}
	\item Maintained DPDK drivers for Intel ixgbe, i40e, and ice NIC families, diagnosing and fixing 12 driver defects
	\item Implemented Header Split on RX path to eliminate per-packet CPU copy overhead
	\item Developed RTP protocol splitting (Intel E810) for direct NIC-to-GPU video data transfer
\end{itemize}
\end{onehalfspacing}

\datedsubsection{\textbf{Baicells Technologies}, Chengdu}{Aug. 2019 -- Nov. 2021}
\role{Linux Kernel Development Engineer}{C kernel MPTCP Userspace IPC}
\begin{onehalfspacing}
CPE Terminal Kernel Module and Userspace Management Development
\begin{itemize}
	\item Developed MPTCP (Multipath TCP) kernel module enabling simultaneous transmission over multiple links (4G+Wi-Fi) for increased bandwidth and reliability; performed in-depth kernel stack RX path performance analysis and optimization
	\item Designed userspace configuration management tool with IPC mechanism for parameter delivery, validation, and runtime state synchronization with kernel module
\end{itemize}
\end{onehalfspacing}

\datedsubsection{\textbf{State Grid Corporation of China}}{Jul. 2013 -- Jul. 2019}
\role{Network Engineer}{OSPF BGP VLAN Firewall Network Architecture}
\begin{onehalfspacing}
Provincial-Level Power Communication Network Planning and Operations
\begin{itemize}
	\item Responsible for network architecture design and planning across 50+ substations, with deep understanding of OSPF/BGP routing protocols, VLAN isolation, and firewall security policies
\end{itemize}
\end{onehalfspacing}

\section{\faGraduationCap\ Education}
\datedsubsection{\textbf{Sichuan University}, Chengdu, China}{2010 -- 2013}
\textit{M.S.} in Software Engineering
\datedsubsection{\textbf{Sichuan University}, Chengdu, China}{2006 -- 2010}
\textit{B.S.} in Software Engineering

\section{\faWrench\ Technical Skills}
\begin{itemize}[parsep=0.5ex]
  \item C Bash Linux RISC-V DPDK SPDK Ceph gdb vim vscode vibe-coding (Claude Code, Cursor)
  \item Network protocol stacks, driver development, distributed storage, kernel/userspace IPC, computer architecture
  \item Git
\end{itemize}

\section{\faGlobe\ Miscellaneous}
\begin{itemize}[parsep=0.5ex]
  \item Languages: Mandarin (Native), English (Proficient -- TOEFL 92, CET-6)
\end{itemize}

\end{document}
